\documentclass[12pt,a4paper]{article}

\usepackage{amsmath,amssymb,amsthm}
\usepackage{hyperref}
\usepackage{doi}
\usepackage[numbers,sort&compress]{natbib}
\usepackage{microtype}
\usepackage{geometry}
\geometry{margin=2.5cm}

\newtheorem{theorem}{Theorem}[section]
\newtheorem{lemma}[theorem]{Lemma}
\newtheorem{corollary}[theorem]{Corollary}
\newtheorem{proposition}[theorem]{Proposition}
\newtheorem{remark}[theorem]{Remark}
\newtheorem{definition}[theorem]{Definition}

\DeclareMathOperator{\Sym}{Sym}
\DeclareMathOperator{\End}{End}
\newcommand{\Heff}{H_{\mathrm{eff}}}
\newcommand{\Leff}{L_{\mathrm{eff}}}
\newcommand{\su}{\mathfrak{su}}
\newcommand{\Heis}{\mathrm{Heis}}
\newcommand{\SU}{\mathrm{SU}}

\title{Temporal Casimir Rigidity and Closure of the Effective Metric:\\
Identification of $A_\tau = 2$\\[6pt]
{\large Q11 of the Cosmochrony Spectral Geometry Programme}}

\author{J\'er\^ome Beau\\
\small Independent Researcher\\
\small \texttt{jerome.beau@cosmochrony.org}}

\date{2026}

\begin{document}

  \maketitle

  \begin{abstract}
    Paper Q5b~\cite{Beau2026q5b} leaves the temporal coefficient $A_\tau$ of the effective co-metric
    $g^{\mu\nu} = \mathrm{diag}(-A_\tau, 2, 2, 2)$ as the sole remaining free parameter of the
    Cosmochrony spectral geometry programme (open problem Q5b-O3).
    The present paper closes this problem by proving $A_\tau = 2$ under the Q5a
    hypotheses~\cite{Beau2026q5a} and the spectral universality hypothesis~\cite{Beau2026u1}~[U].
    The argument proceeds in three steps.
    First, we establish a \emph{cascade internalization lemma}: the temporal coordinate $\tau$,
    identified with the BFS depth index $n$ in Q5b Remark~3.3~\cite{Beau2026q5b}, has $\partial_\tau$
    as the continuum limit of the cascade increment operator $T\colon \sigma_c(n)\mapsto\sigma_c(n+1)$
    by the Carnot--Carath\'eodory convergence of Q5b Theorem~3.2 (Pansu~\cite{Pansu1989}).
    This identification requires no new hypothesis.
    Second, the O-series spectral universality [U], proved in U1~\cite{Beau2026u1}, implies that
    $T$ is asymptotically $\SU(2)$-equivariant, so $\partial_\tau$ acts on
    $\Sym^2(V_\rho)$ without introducing any new irreducible representation.
    Third, Schur's lemma forces the unique $\SU(2)$-invariant quadratic form on $\Sym^2(V_\rho)$
    to be proportional to the Casimir $\mathcal{C}_{\su(2)} = 2\cdot\mathrm{Id}$,
    and Q8~\cite{BeauQ8} fixes the normalisation scalar to unity.
    The effective co-metric is therefore fully determined:
    $g^{\mu\nu} = 2\,\eta^{\mu\nu}$.
    No free parameter remains.
    The temporal direction is not an additional geometric axis; it is the continuum
    infinitesimalisation of the admissible cascade itself.
  \end{abstract}

  \textbf{Keywords:} Heisenberg group; Weil representation; cascade generator; Carnot convergence;
  Casimir operator; $\su(2)$-module; effective metric; temporal ordering; Cosmochrony.

  \textbf{MSC 2020.} 22E25, 81R05, 53B30, 35H10, 22E70.

  \tableofcontents

  \newpage
  \section{Introduction}
    \label{sec:intro}

    \subsection{The remaining open problem}

      The chain of results Q5b--Q10 of the Cosmochrony spectral geometry programme establishes
      an effective Lorentzian co-metric on $\mathbb{R}_\tau\times\Heis_3(\mathbb{R})$ of the form
      \begin{equation}
        g^{\mu\nu} = \mathrm{diag}(-A_\tau,\, 2,\, 2,\, 2),
        \label{eq:metric-intro}
      \end{equation}
      where all three spatial coefficients equal $2$, forced by the $\su(2)$-Casimir
      $\mathcal{C}_{\su(2)}|_{\Sym^2(V_\rho)} = 2\cdot\mathrm{Id}$~\cite{BeauQ8,Beau2026q10}.
      The single remaining unknown is the temporal coefficient $A_\tau > 0$,
      which Q5b~\cite{Beau2026q5b} declares open as problem Q5b-O3: \emph{identify $A_\tau$ from
  the projective temporal ordering data}~\cite{Beau2026pto}.
      The present paper resolves Q5b-O3 by proving $A_\tau = 2$,
      thereby closing $g^{\mu\nu} = 2\,\eta^{\mu\nu}$ with no free parameter.

    \subsection{Main result}

      \begin{theorem}[Temporal Casimir Rigidity]
        \label{thm:main}
        Under the Q5a hypotheses \textup{[H1], [H2], [H-w], [H-E1], [C]}~\cite{Beau2026q5a}
        and the spectral universality hypothesis \textup{[U]} proved in
        U1~\cite{Beau2026u1}, the temporal coefficient of the effective co-metric satisfies
        \[
          A_\tau = 2 = \mathcal{C}_{\su(2)}\big|_{\Sym^2(V_\rho)}.
        \]
        Consequently the effective co-metric on $\mathbb{R}_\tau\times\Heis_3(\mathbb{R})$ is
        \[
          g^{\mu\nu} = 2\,\eta^{\mu\nu},
        \]
        with no remaining free parameter.
        Open problem Q5b-O3 of~\cite{Beau2026q5b} is closed.
      \end{theorem}

      The result is conditional on the same hypotheses as Q10~\cite{Beau2026q10}: no new assumption
      is required beyond what is already in force.

    \subsection{The key observation}

      The proof rests on a single structural observation that was implicit in Q5b but not stated
      explicitly.
      Q5b Remark~3.3~\cite{Beau2026q5b} identifies $\tau$ with the BFS depth index $n$
      (the projective ordering parameter from~\cite{Beau2026pto}).
      By the Carnot--Carath\'eodory convergence theorem of Pansu~\cite{Pansu1989}
      (Q5b Theorem~3.2~\cite{Beau2026q5b}), the discrete cascade increment
      $T\colon\sigma_c(n)\mapsto\sigma_c(n+1)$ passes to the continuum derivative $\partial_\tau$
      in the large-$q$ limit.
      Consequently:
      \begin{center}
        \emph{$\partial_\tau$ is not an external geometric direction; it is the continuum
        infinitesimalisation of the admissible cascade.}
      \end{center}
      This means $\partial_\tau$ acts on the same spectral space $\Sym^2(V_\rho)$ as the
      horizontal directions, and the $\su(2)$-equivariance of $T$ (a consequence of [U]) forces
      $\partial_\tau$ into the same representation-theoretic rigidity as the spatial coefficients.
      Schur's lemma then eliminates $A_\tau$ as a free parameter.

    \subsection{Organisation}

      Section~\ref{sec:setup} fixes notation and recalls the relevant results of Q5b, Q7--Q10,
      and U1.
      Section~\ref{sec:internalization} establishes the cascade internalization lemma (Lemma~\ref{lem:internalization}).
      Section~\ref{sec:equivariance} derives the asymptotic $\SU(2)$-equivariance of the cascade
      generator from [U].
      Section~\ref{sec:schur} applies Schur's lemma and completes the proof of Theorem~\ref{thm:main}.
      Section~\ref{sec:consequences} draws the programme-level consequences.
      Section~\ref{sec:summary} summarises the result and the remaining landscape.

      \newpage
  \section{Setup and Notation}
    \label{sec:setup}

    \subsection{Objects and standing hypotheses}

      Fix a prime $q\geq 5$.
      The discrete Heisenberg group is $G_q = \Heis_3(\mathbb{Z}/q\mathbb{Z})$ with generators
      $X_q,Y_q,Z_q$.
      The Weil representation at central character $e^{2\pi i c/q}$ is
      $\rho_{q,c}\colon G_q\to\mathrm{U}(\mathbb{C}^q)$.
      The BFS fingerprint energy at depth $n$ and character $c$ is
      \[
        \sigma_c(n) := \frac{\delta_{r_n}(c)}{|S_n|},
      \]
      where $r_n$ is the BFS radius at step $n$ and $|S_n|$ is the shell size.
      The effective Hilbert space is $\Heff = \Sym^2(V_\rho)$, the spin-1 $\su(2)$-module
      identified in Q7~\cite{Beau2026q7}.

      The Q5a hypotheses [H1], [H2], [H-w], [H-E1], [C] govern the admissibility filter $\Pi_q$
      and the Mosco convergence of the admissibility forms; see~\cite{Beau2026q5a} for precise
      statements.
      Hypothesis [U] (Q10 Definition~3.1~\cite{Beau2026q10}) asserts the existence of a universal
      profile $\sigma^*(n)$ such that
      \[
        \max_{c\in(\mathbb{Z}/q\mathbb{Z})^\times}\,|\sigma_c(n) - \sigma^*(n)|
        \;\leq\; \varepsilon(q)\,\sigma^*(n),
      \]
      uniformly for $n\leq n^*(q)$, with $\varepsilon(q)=O(q^{-1/2})$.
      Hypothesis [U] is proved in U1~\cite{Beau2026u1}.

    \subsection{Results of Q5b--Q10 in force}

      The following are used without further proof.

      \begin{itemize}
        \item \textbf{Q5b Remark~3.3}~\cite{Beau2026q5b}: the coordinate $\tau$ on $\mathbb{R}_\tau$
        is the BFS depth index $n$, i.e.\ the projective ordering parameter of~\cite{Beau2026pto}.
        \item \textbf{Q5b Theorem~3.2}~\cite{Beau2026q5b}: the rescaled BFS metric space
        $(G_q, n^{-1}d_w)$ converges to the Carnot--Carath\'eodory space
        $(\Heis_3(\mathbb{R}), d_{\mathrm{CC}})$ in the Gromov--Hausdorff sense (Pansu~\cite{Pansu1989}).
        \item \textbf{Q5b Theorem~5.2}~\cite{Beau2026q5b}: the effective operator on
        $\mathbb{R}_\tau\times\Heis_3(\mathbb{R})$ is
        $\Leff = -A_\tau\partial_\tau^2 + A_H\Delta_H$,
        with $A_\tau$ from the projective ordering structure~\cite{Beau2026pto} and
        $A_H$ from the admissibility weight [H-w].
        \item \textbf{Q7}~\cite{Beau2026q7}: the equivariant bridge
        $\phi\colon\Heff\xrightarrow{\sim}W_{\mathrm{sp}}$ is unique up to positive scalar;
        $\Heff = \Sym^2(V_\rho)$ is the irreducible $\su(2)$-module of spin $1$.
        \item \textbf{Q8}~\cite{BeauQ8}: $A_Z = \mathcal{C}_{\su(2)}|_{\Sym^2(V_\rho)} = 2$
        (Casimir rigidity of the central direction).
        \item \textbf{Q10 + U1}~\cite{Beau2026q10,Beau2026u1}: $A_H = 2$ and
        $g^{\mu\nu} = \mathrm{diag}(-A_\tau, 2, 2, 2)$.
      \end{itemize}

  \section{Cascade Internalization}
    \label{sec:internalization}

    \subsection{The cascade increment operator}

      The cascade is the map $U_{n+1} = \Pi\circ\sigma(U_n)$ acting on spectral data.
      At each step $n$, the BFS fingerprint energy $\sigma_c(n)$ is updated.
      Define the \emph{cascade increment operator}
      \[
        T\colon H_{\mathrm{eff}}\longrightarrow H_{\mathrm{eff}},
        \qquad T\,\sigma_c(n) := \sigma_c(n+1).
      \]
      The operator $T$ acts on the same spectral space as the spatial observables:
      its domain and codomain are both $\Heff = \Sym^2(V_\rho)$.

      \begin{lemma}[Cascade Internalization]
        \label{lem:internalization}
        Under Q5b Remark~3.3 and Theorem~3.2~\cite{Beau2026q5b}, the formal derivative $\partial_\tau$
        appearing in the effective operator $\Leff$ is the continuum limit of the cascade increment
        operator $T$.
        More precisely, in the Gromov--Hausdorff limit $q\to\infty$,
        \[
          T - \mathrm{Id} \;\longrightarrow\; \partial_\tau
        \]
        as a derivation on the spectral data, and $\partial_\tau$ acts on $\Sym^2(V_\rho)$.
      \end{lemma}

      \begin{proof}
        By Q5b Remark~3.3~\cite{Beau2026q5b}, $\tau$ is identified with the BFS depth index $n$.
        The continuum limit is established by Q5b Theorem~3.2~\cite{Beau2026q5b} (following
        Pansu~\cite{Pansu1989}): in the limit $q\to\infty$ with $n = \lfloor r q^\alpha\rfloor$
        for $\alpha < 1$, the discrete BFS metric space converges to the Carnot--Carath\'eodory
        space $(\Heis_3(\mathbb{R}), d_{\mathrm{CC}})$ in the Gromov--Hausdorff sense.
        Under this convergence, discrete differences $f(n+1) - f(n)$ pass to the continuum derivative
        $\partial_\tau f(\tau)$ for any sequence $f(n)$ that is Lipschitz in the BFS metric.
        The fingerprint energies $\sigma_c(n)$ are Lipschitz in $n$ within the pre-saturation window
        $n\leq n^*(q)$ (Q5a Hypothesis [H-E1]~\cite{Beau2026q5a}).
        Since $T\,\sigma_c(n) = \sigma_c(n+1)$ and the domain of $T$ is $\Heff = \Sym^2(V_\rho)$,
        the derivation $\partial_\tau$ inherits the same domain.
      \end{proof}

      \begin{remark}
        Lemma~\ref{lem:internalization} requires no hypothesis beyond Q5b and the Q5a Lipschitz
        condition [H-E1].
        It makes explicit an identification that was implicit in Q5b but left unstated:
        the temporal derivative $\partial_\tau$ is not an external geometric direction imposed on
        $\mathbb{R}_\tau$, but the infinitesimal generator of the admissible cascade, already
        acting on $\Sym^2(V_\rho)$.
      \end{remark}

    \subsection{No new irreducible representation}

      A direct consequence of Lemma~\ref{lem:internalization} is the following.

      \begin{corollary}
        \label{cor:no-new-rep}
        As an operator on $\Sym^2(V_\rho)$, the cascade generator $\partial_\tau$ lies in
        $\End(\Sym^2(V_\rho))$.
        The tensor product $\mathbb{R}_\tau\otimes\Sym^2(V_\rho)$ decomposes as
        $\mathbf{1}\otimes\Sym^2(V_\rho) \cong \Sym^2(V_\rho)$ as $\SU(2)$-modules,
        where $\mathbf{1}$ denotes the trivial one-dimensional representation.
        In particular, the temporal direction introduces no new irreducible representation of $\SU(2)$.
      \end{corollary}

      \begin{proof}
        The cascade increment $T - \mathrm{Id}$ acts on $\Heff = \Sym^2(V_\rho)$ by definition.
        Its continuum limit $\partial_\tau$ (Lemma~\ref{lem:internalization}) therefore acts on the same
        module.
        The continuous coordinate $\tau$ is a scalar parameter, so the $\tau$-sector carries the trivial
        $\SU(2)$-representation $\mathbf{1}$.
        Tensoring by $\mathbf{1}$ does not alter the $\SU(2)$-module structure:
        $\mathbf{1}\otimes\Sym^2(V_\rho)\cong\Sym^2(V_\rho)$.
      \end{proof}

  \section{Asymptotic \texorpdfstring{$\SU(2)$}{SU(2)}-Equivariance of the Cascade Generator}
    \label{sec:equivariance}

    The following proposition reformulates [U] as a statement about representation-theoretic
    equivariance of the cascade, replacing the analytic uniformity condition by a categorical one.

    \begin{proposition}[$\SU(2)$-Equivariance of $T$]
      \label{prop:equivariance}
      Under the Q5a hypotheses and \textup{[U]} (proved in U1~\cite{Beau2026u1}), the cascade
      increment operator $T$ satisfies
      \[
        T\circ\rho(g) = \rho(g)\circ T + O(q^{-1/2}),
        \qquad\forall\,g\in\SU(2),
      \]
      as operators on $\Sym^2(V_\rho)$, uniformly for $n\leq n^*(q)$.
      In the limit $q\to\infty$, $T$ is $\SU(2)$-equivariant.
    \end{proposition}

    \begin{proof}
      Hypothesis [U] asserts
      $\max_{c}|\sigma_c(n) - \sigma^*(n)| \leq \varepsilon(q)\,\sigma^*(n)$
      with $\varepsilon(q) = O(q^{-1/2})$.
      The increment satisfies
      \[
        T\,\sigma_c(n) - \sigma^*(n+1)
        = \sigma_c(n+1) - \sigma^*(n+1)
        \;\leq\; \varepsilon(q)\,\sigma^*(n+1).
      \]
      Since $\sigma^*(n)$ is $c$-independent, it is invariant under the $\SU(2)$-action on characters
      (which permutes $c\in(\mathbb{Z}/q\mathbb{Z})^\times$).
      For any $g\in\SU(2)$ acting as a permutation $\pi_g$ on characters,
      \begin{align*}
        \bigl\|T(\rho(g)\,\sigma_c(n)) - \rho(g)(T\,\sigma_c(n))\bigr\|
        &= \bigl\|\sigma_{\pi_g(c)}(n+1) - \sigma_{\pi_g(c)}(n+1)\bigr\|
        \;\leq\; 2\varepsilon(q)\,\sigma^*(n+1).
      \end{align*}
      As $\varepsilon(q)\to 0$, this bound vanishes, giving asymptotic equivariance.
    \end{proof}

    \begin{remark}
      Proposition~\ref{prop:equivariance} is the representation-theoretic reading of [U].
      Analytically, [U] asserts uniform $c$-independence of the spectral profiles.
      Categorically, it asserts that the cascade morphism $T$ commutes with the $\SU(2)$-action
      in the limit.
      These are equivalent statements; the categorical reading is what enters the Schur argument.
    \end{remark}

  \section{Schur Rigidity and Proof of Theorem~\ref{thm:main}}
    \label{sec:schur}

    \subsection{The Schur argument}

      \begin{proposition}[Uniqueness of the invariant quadratic form]
        \label{prop:schur}
        Let $V = \Sym^2(V_\rho)$ be the irreducible $\su(2)$-module of spin $1$.
        Every $\SU(2)$-equivariant quadratic form $Q\colon V\to\mathbb{R}$ is of the form
        \[
          Q = \alpha\,\mathcal{C}_{\su(2)}\big|_V,
          \qquad\alpha\in\mathbb{R}.
        \]
      \end{proposition}

      \begin{proof}
        By Schur's lemma, every $G$-equivariant endomorphism of an irreducible representation is
        a scalar multiple of the identity.
        The quadratic form $Q$ defines a $\SU(2)$-equivariant symmetric bilinear form on $V$,
        hence an equivariant map $V\to V^*\cong V$ (the isomorphism using the Killing form).
        Since $V = \Sym^2(V_\rho)$ is irreducible of spin $1$ (Q7~\cite{Beau2026q7}),
        this map is a scalar, say $\alpha\,\mathrm{Id}$.
        The unique normalisation-invariant scalar is the Casimir eigenvalue
        $\mathcal{C}_{\su(2)}|_V = 2$ (Q8~\cite{BeauQ8}).
      \end{proof}

    \subsection{Proof of Theorem~\ref{thm:main}}

      \begin{proof}[Proof of Theorem~\ref{thm:main}]
        The effective co-metric coefficient $A_\tau$ is the coefficient of $k_\tau^2$ in the principal
        symbol $\sigma_2(\Leff)$ of the effective operator (Q5b Theorem~5.2~\cite{Beau2026q5b}).
        It is therefore the normalisation scalar of the quadratic form on the $\tau$-sector of
        $T^*(\mathbb{R}_\tau\times\Heis_3(\mathbb{R}))$.

        \emph{Step~1 (Internalization).}
        By Lemma~\ref{lem:internalization}, $\partial_\tau$ is the continuum limit of $T$ acting on
        $\Heff = \Sym^2(V_\rho)$.
        By Corollary~\ref{cor:no-new-rep}, the $\tau$-sector carries the trivial $\SU(2)$-representation
        $\mathbf{1}$, and the full $\tau$-module is $\mathbf{1}\otimes\Sym^2(V_\rho)\cong\Sym^2(V_\rho)$.

        \emph{Step~2 (Equivariance).}
        By Proposition~\ref{prop:equivariance}, $T$ is asymptotically $\SU(2)$-equivariant under [U].
        In the limit $q\to\infty$, the quadratic form defined by $A_\tau k_\tau^2$ on the $\tau$-sector
        must be $\SU(2)$-invariant.

        \emph{Step~3 (Schur).}
        By Proposition~\ref{prop:schur}, every $\SU(2)$-invariant quadratic form on $\Sym^2(V_\rho)$
        is $\alpha\,\mathcal{C}_{\su(2)}|_{\Sym^2(V_\rho)}$ for some scalar $\alpha$.

        \emph{Step~4 (Normalisation).}
        The bridge $\phi\colon\Heff\xrightarrow{\sim}W_{\mathrm{sp}}$ of Q7~\cite{Beau2026q7} is unique
        up to positive scalar, with normalisation fixed by the Casimir eigenvalue
        $\mathcal{C}_{\su(2)}|_{\Sym^2(V_\rho)} = 2$ (Q8~\cite{BeauQ8}).
        Since $\partial_\tau$ is the continuum limit of the same cascade $T$ that enters the spatial bridge
        (Step~1), the scalar $\alpha$ is the same in both sectors.
        Q8 establishes $\alpha = 1$ for the spatial sectors (giving $A_Z = A_H = 2$).
        Therefore $\alpha = 1$ in the $\tau$-sector as well, and $A_\tau = \mathcal{C}_{\su(2)} = 2$.
      \end{proof}

  \section{Consequences}
    \label{sec:consequences}

    \subsection{Closure of the effective metric}

      \begin{corollary}[Isotropic Lorentzian metric]
        \label{cor:metric}
        Under the Q5a hypotheses and \textup{[U]}, the effective co-metric on
        $\mathbb{R}_\tau\times\Heis_3(\mathbb{R})$ is
        \[
          g^{\mu\nu} = 2\,\eta^{\mu\nu},
          \qquad\eta = \mathrm{diag}(-1,+1,+1,+1),
        \]
        of Lorentzian signature $(-,+,+,+)$, with no free parameter.
        After absorbing the conformal factor $\Omega^2 = 2$ into the normalisation of the propagation
        speed $c$ (replacing $c$ by $c/\sqrt{2}$), the effective metric is exactly Minkowski.
      \end{corollary}

      \begin{proof}
        Combine Theorem~\ref{thm:main} ($A_\tau = 2$) with Q10~\cite{Beau2026q10} ($A_H = 2$) and
        Q8~\cite{BeauQ8} ($A_Z = 2$).
        The Lorentzian signature follows from Q5b Theorem~6.1~\cite{Beau2026q5b}.
      \end{proof}

      \begin{remark}[Structural significance of the conformal factor $2$]
        The value $2$ is not a choice of units.
        It is forced by the eigenvalue of the $\su(2)$-Casimir on the spin-$1$ module
        $\Sym^2(V_\rho)$, which is the unique irreducible representation compatible with the
        admissibility constraints (Q7~\cite{Beau2026q7}).
        The same integer governs all four metric components: two horizontal Heisenberg directions
        (Q10), the central Heisenberg direction (Q8), and the temporal direction (this paper).
        This uniform emergence from a single representation-theoretic datum is the structural
        content of the metric closure.
      \end{remark}

    \subsection{Closure of Q5b-O3}

      \begin{corollary}[Resolution of Q5b-O3]
        Open problem Q5b-O3 of~\cite{Beau2026q5b}---\emph{identify $A_\tau$ from the projective
        temporal ordering data}---is closed: $A_\tau = 2$, derived from the cascade internalization
        lemma, [U], and Q8.
      \end{corollary}

    \subsection{The complete derivation chain}

      The closure of the effective metric is now traced through the following chain,
      where each arrow represents a proved theorem or lemma:
      \begin{center}
        \begin{tabular}{cl}
          Foundation~\cite{Beau2026m}       & time $=$ projective irréversible ordering \\
          $\downarrow$                       & \\
          Q5b Remark~3.3~\cite{Beau2026q5b} & $\tau \leftrightarrow n$ (BFS depth) \\
          $\downarrow$                       & \\
          Lemma~\ref{lem:internalization}    & $\partial_\tau \leftrightarrow T$ (cascade generator) \\
          $\downarrow$                       & \\
          [U] + U1~\cite{Beau2026u1}         & $T$ is $\SU(2)$-equivariant \\
          $\downarrow$                       & \\
          Proposition~\ref{prop:schur}       & unique $\SU(2)$-invariant form on $\Sym^2(V_\rho)$ \\
          $\downarrow$                       & \\
          Q8~\cite{BeauQ8}                   & Casimir normalisation $= 2$ \\
          $\downarrow$                       & \\
          Theorem~\ref{thm:main}             & $g^{\mu\nu} = 2\,\eta^{\mu\nu}$.
        \end{tabular}
      \end{center}

    \subsection{Programme-level closure}

      With $A_\tau = 2$ established, the effective geometry programme Q5a--Q11 is complete:

      \begin{itemize}
        \item \textbf{Q5a}~\cite{Beau2026q5a}: admissibility weight $A$; Q5a hypotheses
        [H1], [H2], [H-w], [H-E1], [C].
        \item \textbf{Q5b}~\cite{Beau2026q5b}: BFS--Carnot convergence; effective operator
        $\Leff$ on $\mathbb{R}_\tau\times\Heis_3(\mathbb{R})$; Lorentzian signature.
        \item \textbf{Q7}~\cite{Beau2026q7}: bridge $\phi\colon\Heff\xrightarrow{\sim}W_{\mathrm{sp}}$;
        $\SU(2)$ rigidity.
        \item \textbf{Q8}~\cite{BeauQ8}: $A_Z = 2$ (Casimir rigidity, central direction).
        \item \textbf{Q9}~\cite{Beau2026q9}: [H-lift] proved; lifting is unconditional.
        \item \textbf{Q10}~\cite{Beau2026q10}: $A_H = 2$ (asymptotic spatial isotropy, [U]).
        \item \textbf{U1}~\cite{Beau2026u1}: [U] proved with rate $O(q^{-1/2})$.
        \item \textbf{Q11 (this paper)}: $A_\tau = 2$; $g^{\mu\nu} = 2\,\eta^{\mu\nu}$; Q5b-O3 closed.
      \end{itemize}

  \section{Summary and Open Problems}
    \label{sec:summary}

    \subsection{What this paper establishes}

      \begin{enumerate}
        \item The temporal coordinate $\tau$ is the BFS depth index $n$ (Q5b Remark~3.3), and the
        formal derivative $\partial_\tau$ is the continuum limit of the cascade increment operator
        $T$ (Lemma~\ref{lem:internalization}).
        This requires no hypothesis beyond Q5b and [H-E1].
        \item Under [U] (U1~\cite{Beau2026u1}), the cascade generator $T$ is asymptotically
        $\SU(2)$-equivariant (Proposition~\ref{prop:equivariance}), so $\partial_\tau$ acts on
        $\Sym^2(V_\rho)$ without introducing a new irreducible.
        \item Schur's lemma (Proposition~\ref{prop:schur}) and the Casimir normalisation of Q8
        (Theorem~\ref{thm:main}) force $A_\tau = 2$.
        \item The effective co-metric is $g^{\mu\nu} = 2\,\eta^{\mu\nu}$ with no free parameter
        (Corollary~\ref{cor:metric}).
        Q5b-O3 is closed.
      \end{enumerate}

    \subsection{Open problems}

      \begin{description}
        \item[Q11-O1 (Finite-$q$ cascade equivariance).]
        Proposition~\ref{prop:equivariance} establishes $\SU(2)$-equivariance of $T$ only
        asymptotically as $q\to\infty$.
        An explicit finite-$q$ bound on the equivariance error, beyond the $O(q^{-1/2})$ estimate
        inherited from [U], would sharpen the result.
        \item[Q11-O2 (Unconditional closure).]
        The identification $A_\tau = 2$ is conditional on [U].
        Since [U] is proved in U1~\cite{Beau2026u1} under the Q5a hypotheses alone, the result
        is unconditional modulo the Q5a hypotheses.
        An independent derivation of $A_\tau$ from the Born--Infeld saturation data of~\cite{Beau2026pto}
        would provide a complementary confirmation.
        \item[Q11-O3 (Physical interpretation of the conformal factor).]
        The conformal factor $\Omega^2 = 2$ absorbed into $c$ has a representation-theoretic origin
        (Casimir eigenvalue on the spin-$1$ module).
        Its possible observational consequences---for instance, as a correction to the effective
        speed of light at finite $q$---remain to be explored.
      \end{description}

  \section*{Acknowledgements}

    The author acknowledges the use of large language models as a supportive tool for refining
    language, structure, and internal consistency during the development of this manuscript.

    \newpage
    \bibliography{cosmochrony-bibliography}
    \bibliographystyle{plainnat}

\end{document}
